\documentclass[11pt,a4paper]{article}
\usepackage[T1]{fontenc}
\usepackage[utf8]{inputenc}
\usepackage{lmodern}
\usepackage{amsmath,amssymb,amsthm,mathtools}
\usepackage{geometry}
\geometry{margin=2.6cm}
\usepackage{booktabs,array,tabularx}
\usepackage{enumitem}
\usepackage{fvextra}
\usepackage[dvipsnames]{xcolor}
\usepackage[hidelinks]{hyperref}
\usepackage{xurl}
\hypersetup{pdftitle={Two computations of the height inequality and the multiradial algorithm},pdfauthor={Toshihisa Urahigashi}}
\setlist{itemsep=2pt,topsep=4pt}
\raggedbottom
\emergencystretch=3em
\fvset{fontsize=\fontsize{7}{8.4}\selectfont,frame=single,framesep=1.5mm,breaklines=true}

\newcommand{\Z}{\mathbb{Z}}
\newcommand{\Q}{\mathbb{Q}}
\newcommand{\Ocal}{\mathcal{O}}
\newcommand{\Vbad}{\mathbb{V}_{\mathrm{mod}}^{\mathrm{bad}}}
\newcommand{\qq}{\underline{\underline{q}}}
\newcommand{\TT}{\underline{\underline{\Theta}}}
\newcommand{\leanus}{\textunderscore\allowbreak}
\newcommand{\lean}[1]{{\ttfamily\let\_\leanus #1}}
\newcommand{\orc}[1]{\texttt{[#1]}}

% labels
\newcommand{\lab}[2]{{\upshape\fcolorbox{#1}{#1!8}{\scriptsize\sffamily\bfseries\textcolor{#1}{#2}}}}
\newcommand{\tagS}{\lab{blue}{S}}
\newcommand{\tagX}{\lab{Purple}{X}}
\newcommand{\tagP}{\lab{ForestGreen}{P}}
\newcommand{\tagH}{\lab{BurntOrange}{H}}
\newcommand{\q}[1]{``#1''}
\newcommand{\srcquote}[2]{\begin{quote}\small #1\par\nopagebreak\hfill{\footnotesize #2}\end{quote}}

% written by make_generated.py from ../c, ../cpp and ../outputs; do not edit
\newcommand{\thedoi}{10.5281/zenodo.23968299}
\newcommand{\cppitems}{204}
\newcommand{\cppfiles}{16}
\newcommand{\cpplines}{1{,}170}
\newcommand{\selftestnaive}{-5.1\times10^{18}}
\newcommand{\scanCvalues}{1247.14 and 1235.66}
\newcommand{\tabletwo}{%
D211 & $305.21$ & $77.68$ & $-21.801$ & $135.26$ & $1145.40$ & $79.94$ & $1247.14$\\
D109 & $267.60$ & $68.12$ & $-19.114$ & $116.65$ & $1098.32$ & $70.09$ & $1187.52$\\
D2 & $240.09$ & $76.92$ & $-17.149$ & $158.00$ & $1155.16$ & $62.88$ & $1235.19$\\
D3 & $1128.24$ & $270.90$ & $-80.589$ & $449.24$ & $1871.53$ & $295.49$ & $2247.61$\\
D18 & $614.95$ & $160.86$ & $-43.925$ & $285.92$ & $1478.09$ & $161.06$ & $1683.07$\\
D15 & $516.07$ & $137.29$ & $-36.862$ & $243.31$ & $1384.66$ & $135.16$ & $1556.68$\\
D2357$^\ast$ & $3717.84$ & $30.23$ & $-11.840$ & $-22675.23$ & $-8246.15$ & $24463.93$ & $16229.62$\\}
\newcommand{\tablesix}{%
D211 & $N\le2$ (``no'' from $N=3$) & $N\le15$ (fails from $N=16$)\\
D109 & $N\le2$ (``no'' from $N=3$) & $N\le16$ (fails from $N=17$)\\
D2 & $N\le3$ (``no'' from $N=4$) & $N\le19$ (fails from $N=20$)\\
D3 & $N\le2$ (``no'' from $N=3$) & $N\le6$ (fails from $N=8$)\\
D18 & $N\le2$ (``no'' from $N=3$) & $N\le10$ (fails from $N=11$)\\
D15 & $N\le2$ (``no'' from $N=3$) & $N\le11$ (fails from $N=12$)\\}
\newcommand{\iutcore}{\VerbatimInput{excerpt_iut_core.c}}
\newcommand{\cppbox}{\VerbatimInput{excerpt_box.hpp}}
\newcommand{\cpprigidity}{\VerbatimInput{excerpt_rigidity.cpp}}
\newcommand{\cppstepxi}{\VerbatimInput{excerpt_stepxi.cpp}}
\newcommand{\datafile}{\VerbatimInput{excerpt_data_D211.txt}}


\title{Two computations of the height inequality\\ and the multiradial algorithm:\\ reading IUT from the viewpoint
of computer science}
\author{Toshihisa Urahigashi\\\small ORCID 0009-0004-0460-6242\\\small\texttt{123026.urahigashi.toshihisa@gmail.com}}
\date{Version 1.1.0, 8 November 2026\\\small\href{https://doi.org/\thedoi}{doi:\thedoi}}

\begin{document}
\maketitle

\begin{abstract}
This note contains almost no new mathematical results. It is an essay that proposes a viewpoint and an idea: to read
inter-universal Teichm\"uller theory (IUT) from the side of computer science. It presents the tools for doing so ---
two programs that read the same input, and a mockup of the multiradial algorithm in which some parts are left as
oracles --- and a little analysis made with them. It proves no conclusion. The view that IUT computes the same
quantity in two ways is Mochizuki's own and is also recorded in the interim report of Project LANA. The explicit
computation of the IUT side is printed in the proof of Theorem~1.10 of IUT~IV and redone by Dupuy--Hilado and by Zhou.
Earlier, the toolkit motizukipdf of D.~Richards (October 2026) turned the objects and processes of IUT into C data
types, a web application and a Lean port, with a numerical form of Step~(xi), and the Lean repository
Takkun-kohinata/IUT\_LEAN (July 2026) proved, for one reading, that the containment of Step~(v) together with
Corollary~3.12 gives the height inequality, and conversely. What this note adds is to place the classical computation
and the computation of IUT side by side as programs that read the same input, the mockup with oracles separated by
types, and a computer-science reading together with a position (\S1.3).

The argument has five points. (1) In each of the 36 readings of \cite{RF}, whose axes come from the author's notes and
in part from lana-agents/iut, Corollary~3.12 of IUT~III and Theorem~1.10 of IUT~IV together give the same height
inequality $(k-1)|\log\qq|\le C$ \cite{RF}. (2) This height inequality is written in classical quantities only; for a given curve,
both sides can be computed without IUT. (3) What IUT supplies is a different method of computing the two sides;
Mochizuki himself describes the development of a mechanism for computing the effect on heights as the fundamental
theme of IUT \cite{Alien}. (4) The key is universality. Corollary~3.12 claims to hold for every initial $\Theta$-datum.
Since the value for a single datum can be computed classically, what IUT has to add is the ``for every''. (5) What is
supposed to supply this universality is the multiradial algorithm of Theorem~3.11. This note regards it as the essence
of IUT, and takes the position that one has not understood IUT essentially unless this algorithm has been elucidated
and implemented.

The tools are two. One is a pair of C programs: the classical computation, and the computation of IUT in the reading
of IUT~IV. The other is a mockup of the multiradial algorithm in C++. In the mockup, the parts that can be implemented
at present are implemented: the local invariants from the abstract group $G_{K_v}^{\mathrm{ab}}$, the log-shells,
Kummer classes at a finite level, cyclotomic rigidity from the commutator in a theta group, the labels and the tuples
of the tensor packets, and the hull of the possible images. The rest is left as named oracles: the anabelian
reconstructions, the Kummer classes of the theta values, the global data, the containment of (Ind3), and Step~(xi).
The values of the oracles are reverse-computed from the ring structure of a theater or from the conclusion, and the
output says so. On seven explicit initial $\Theta$-data the mockup agrees with the C programs on \cppitems{} items.
The implemented part is separated by types so that it cannot read the ring structure of a theater, and its output does
not depend on the coordinates that a theater hides. In the family $q\to q^N$, the implemented part runs the same steps
for every $N$, while the reverse-computed Step~(xi) answers ``no'' from $N=3$ (from $N=4$ for one datum) on the six
completely factored data.

lana-agents/iut has shown that Step~(xi) does not follow from the formalized content of Theorem~3.11: over that
content Step~(xi) is equivalent to the Scholze--Stix inequality, and the missing input is of ABC strength \cite{lana}.
This is the result of taking the content of Theorem~3.11, as currently written, as a specification; as lana-agents/iut
itself states, it does not show that Step~(xi) fails for some genuine datum. In the view of this note, the key lies in the
implementation of the multiradial algorithm. The future task is to implement the oracles as actual algorithms, or as
approximations to them. \textbf{No claim is made about the truth of IUT, of Corollary~3.12, or of the abc
conjecture.}
\end{abstract}

\tableofcontents

\section{Introduction}

\subsection{What this note is}
As mathematical work, this note has little content. It contains hardly any theorem or new estimate. What it offers
instead is a viewpoint and tools for trying it out. The viewpoint is to compare the way in which IUT computes the
height inequality with the classical way, \emph{as programs}. If this is done, the tools of computer science ---
program equivalence, information flow, relational verification, mockups with oracles --- may let one approach the
content of IUT.

As the position of this note, two things are stated: ``the essence of IUT is the multiradial algorithm'' and ``one has
not understood IUT essentially unless this algorithm has been elucidated and implemented''. Neither is a proof; both
are statements of a viewpoint (\S6.2).

\subsection{Where this note stands}
This note is also a summary of a series of notes on the computation of the height inequality in IUT. The aim of every
one of them has been to grasp the exact algorithm by which IUT computes the height inequality.

\begin{center}\small
\begin{tabularx}{\textwidth}{@{}l>{\raggedright}p{6.3cm}X@{}}
\toprule
Key & Title (version) & Content\\
\midrule
\cite{First} & An explicit initial $\Theta$-datum with mixed reduction at a split prime (0.1.7) & the initial
$\Theta$-datum $D_{211}$ over $\Q(\sqrt5)$, with one bad and one good place above $211$\\
\cite{Second} & An initial $\Theta$-datum over an imaginary quadratic field, and numerical values of the local quantities
of IUT~IV (1.0.5) & the datum $D_{109}$ over $\Q(i)$ and the local quantities of IV on explicit data\\
\cite{Third} & The cost of each distinction in the local computations of IUT~IV (1.0.2) & what each separation in the
estimates of IV is worth on explicit data\\
\cite{Fourth} & Split bad primes and the tuple structure of IUT (1.0.4) & split bad primes and tuples; (Ind2) acts by
unit scalars; the procession-normalized coefficient\\
\cite{Fifth} & Two constants in the numerical estimates of IUT~IV (1.0.0) & $12/(\ell(\ell+1))$ compared with
$12/\ell^2$\\
\cite{Cat} & Explicit initial $\Theta$-data and the two readings of Step~(v) of IUT~IV, Theorem~1.10: a catalogue (1.0.0) &
seven explicit initial $\Theta$-data and two readings of Step~(v)\\
\cite{Stepv} & A local non-containment in Step~(v) of IUT~IV, Theorem~1.10 (3.0.0) & a local non-containment in
Step~(v), conditional on the readings (M), (R)\\
\cite{SL} & Segment or loop? (1.1.0) & an overview of the dispute over Corollary~3.12 in homotopy type theory\\
\cite{II} & Reading the IUT dispute through indexed identifications (2.0.0) & an educational model in undergraduate
mathematics; readings of (Ind1) and a conservation law (Theorem~14.3)\\
\cite{WT} & A programmer's walkthrough of inter-universal Teichm\"uller theory (1.1.0) & a program that runs the parts of
IUT that can be written as procedures on explicit data\\
\cite{Gap} & The gap between the two walls is positive (1.0.0) & a closed form for the (Ind2)-terms of lana-agents/iut;
$\Delta(D)>0$\\
\cite{RF} & No reading avoids the height inequality (1.0.0) & the same height inequality in each of 36 readings; the
conservation of the load\\
\bottomrule
\end{tabularx}

\smallskip Table 1: the earlier notes (all on Zenodo; the last seven under CC~BY~4.0).
\end{center}

\subsection{Prior work}
Much of the viewpoint of this note has precedents. We state first how much is new (the search for prior work was
done on 5 November 2026).

\tagS{} \emph{The view that there are two computations is Mochizuki's own.} Remark~3.11.1 of IUT~III speaks of the
construction of Theorem~3.11 as an algorithm and calls the theta values its hidden internal structures. It also
distinguishes Corollary~3.12, which concerns only the qualitative logical consequences of the algorithm, from the
explicit computation of IUT~IV, \S1, which uses essentially the way in which theta values occur:
\srcquote{\q{That is to say, theta functions/theta values may be regarded as HIS of the construction algorithm of
Theorem 3.11 --- somewhat like the internal structure of the CPU or operating system of a computer!} \dots{}
\q{On the other hand, we observe that, unlike Corollary 3.12 below, which only concerns qualitative logical
aspects/consequences of the construction algorithm of Theorem 3.11, the explicit computation to be performed in
[IUTchIV], \S1, of the log-volumes that occur in the statement of Corollary 3.12 makes essential use of the way in
which theta values occur in the construction algorithm of Theorem 3.11.}}{IUT~III, Remark~3.11.1, (iv)
\cite{IUT3}}

\tagX{} The interim report of Project LANA records the same view as Mochizuki's:
\srcquote{\q{According to Mochizuki, however, IUT provides two distinct ways of computing the log-volume of the
$q$-pilot object: the usual computation and the computation via the multiradial algorithm.}}{LANA interim report,
\S10.3 \cite{LANA}}

\tagS{} \emph{The explicit computation of the IUT side} is printed in the proof of Theorem~1.10 of IUT~IV itself: in
Step~(v) the average over the tuples of the tensor packets is taken by Proposition~1.7 of IV, and the average over the
labels by the identities (E1), (E2), with the coefficients $\frac{\ell+5}4$ and $\frac{\ell+1}{24}$ (\S10\,(ii))
\cite{IUT4}. \tagX{} Dupuy--Hilado redo it \cite{DH}. They take the average over the
tuples of the tensor packets and over the labels as an ``iterated expectation'', and obtain the average
$\ell(\ell+1)/12$ of $j^2$ over the labels, the average $\frac{\ell+5}4$ of the number of factors, and the Szpiro
constant $6+\varepsilon_\ell$ with $\varepsilon_\ell=(24\ell+72)/(\ell^2+\ell-12)=O(1/\ell)$. Observations (ii)--(iv)
of \S10 are a re-check of this computation in the form of programs. Zhou gives an algorithm computing an estimate of
the same form and has turned it into Python code \cite{Zhou}; the code of Part~II of his series is in C++ and C, with
Jupyter notebooks \cite{ZhouII}.
\srcquote{\q{we can view the expectation of $X$ as an `iterated expectation'\,}}{Dupuy--Hilado, \S7.1 \cite{DH}}
\srcquote{\q{Then there exists an algorithm [whose construction is presented in the proof] to compute a real number
$\mathrm{Vol}(\mathfrak R)\ge0$ that depends only on~$\mathfrak R$}}{Zhou, Corollary~1.13 \cite{Zhou}}

\tagX{} \emph{Separating a finite layer from a structural layer} is in Maeki \cite{Maeki}, who isolates the layer of
coefficients of the explicit estimates of IUT~IV that can be checked finitely, and delegates the structural layer of
IUT (Hodge theaters, Frobenioids, mutually alien copies, log-volume inequalities) to an assumption IUT-Struct. The
mockup of \S9, which names the parts that cannot be computed and leaves them as oracles, divides in the same way.
\srcquote{\q{the IUT-Explicit-ADIC isolates the $\Sigma_1$-verifiable finite content of the explicit IUT
estimates, while delegating the truth of the full structural IUT layer to the IUT-Struct assumption.}}{Maeki, the
abstract of the Zenodo record \cite{Maeki}}

\tagS{} \emph{The meaning of ``algorithm''.} A ``group-theoretic algorithm'' in the sense of IUT is a construction
phrased in a language that depends only on the topological group structure of a fundamental group, not a procedure
that runs on a computer \cite{AbsTopIII}. \tagX{} Joshi objects that it is not a finitistic procedure and that it is
highly ineffective \cite{Joshi}. What this note calls an ``implementation'' (\S6.2) means closing this gap.
\srcquote{\q{We use the term `group-theoretic algorithm' to mean that the algorithm in question is phrased in
language that only depends on the topological group structure of the fundamental group under
consideration.}}{Mochizuki, Topics in absolute anabelian geometry~III, Remark~1.9.8 \cite{AbsTopIII}}
\srcquote{\q{But `algorithm' as used in (Mochizuki, 2021) is not a finitistic procedure as commonly understood by
mathematicians.} \dots{} \q{Furthermore, these procedures are highly ineffective, notably they do not produce the claimed
outputs, and do not certifiably produce outputs which can be distinguished from each other (this is necessary in the
proof).}}{Joshi, comment~(5) \cite{Joshi}}

\tagX{} \emph{Turning the objects and processes of IUT into programs} was done earlier in the toolkit motizukipdf of
D.~Richards (GitHub, created on 27 October 2026; consulted at commit \texttt{2665db1} of 30 October 2026)
\cite{Richards}. It turns the objects of IUT~I--IV into C data types with a registry of processes, runs the C code in
the browser through WebAssembly, and has a Lean port (\lean{lean/IUT/LeanForm.lean}). Its data model lets a
$\Theta$-link share the unit groups but never the ring structures. Its process \texttt{ComputeThmOneTen} evaluates the
constants and the final inequality of Theorem~1.10 of IUT~IV; its process \texttt{Fuzzify} widens the $\Theta$-pilot to
an interval and reports whether the interval contains $-|\log\qq|$, a numerical form of the comparison of Step~(xi); its
Lean port proves APT for its model (\lean{dsc\_APT}) and, from the comparison in the common container taken as a
hypothesis, the bound $\log\mathfrak q\le\frac{12\ell(\ell^*+3)}{\ell(\ell+1)-12}\log\mathfrak d^K$, with
$\ell^*=(\ell-1)/2$ (\lean{corThreeTwelve\_unfolded}), whose coefficient is the constant $c_\ell$ of observation~(iv) of
\S10. Version 1.0.0
of this note did not cite it, although the author's walkthrough~\cite{WT} does.

\tagX{} \emph{Points (1) and (4) of the argument (\S1.4) were made earlier, for one reading,} in the Lean
repository Takkun-kohinata/IUT\_LEAN (GitHub, 15 July 2026) \cite{TK}. There the containment of Step~(v) of IV, with
$\lambda=\mathrm{ord}(q^{j^2})$ read in the normalization in which Corollary~3.12 measures the $q$-pilot, together with
Corollary~3.12 gives $(\ell(\ell+1)/12-1)\,|\log\qq|\le L$ for its ledger $L\ge0$ of IV, Propositions~1.1--1.4
(\lean{thetaHull\_forces\_height}); conversely, this inequality together with the condition $C_\Theta\ge-1$ of IV,
Theorem~1.10 gives a model with that containment (\lean{thetaHullTranscription\_iff}); and with $L=0$ the containment
cannot hold (\lean{no\_lossless\_theta\_hull}). The accompanying request in the same repository states the consequence
for curves:
\srcquote{\q{the satisfiability of (P) for actual curves is logically equivalent to the ledger absorbing the Gauss
gain, i.e.\ to $L\ge(l(l+1)/12-1)\cdot|\log(q)|$ --- the height inequality itself.}}{Takkun-kohinata/IUT\_LEAN,
\lean{REQUEST\_theta\_hull\_derivation.md}, English abstract \cite{TK}; symbols typeset as mathematics}
Its observation (F3), in the Japanese part, adds that the ledger does not depend on $|\log(q)|$, so that supplying the
containment uniformly for a family of curves of unbounded height is equivalent to supplying the height inequality
(paraphrased).

\tagH{} \emph{What this note adds}, in addition to the work above, is limited to three things. (a) The classical
computation, written from the statement of Theorem~1.10 alone, and the computation of IUT in the reading of IV, as two
programs that read the same local data of seven explicit initial $\Theta$-data and agree at the two ends (\S8). (b) A
mockup in which a part of the multiradial algorithm is left as oracles whose values are reverse-computed, with a
separation of information flow by types, a check of non-interference, and a scan of the family $q\to q^N$ (\S9). (c) A
computer-science reading (non-interference, relations over all inputs) and the position that ``one has not understood
IUT essentially unless the multiradial algorithm has been elucidated and implemented''. Mochizuki himself uses a word
close to ``non-interference'' in a mathematical sense (\S6.1); we do not claim that the reading of (c) is the same as
his.

\subsection{The shape of the argument}
\begin{enumerate}
\item \textbf{In every reading, IUT passes through the height inequality} (\S2). A summary of \cite{RF} and of
lana-agents/iut \cite{lana}.
\item \textbf{The height inequality is classical} (\S3). Both sides are written in classical quantities and can be
computed for a given curve without IUT. Whether it holds for all curves is open (a form of the Szpiro conjecture).
\item \textbf{What IUT supplies is a different method of computing the two sides} (\S4). We quote Mochizuki and state
our reading of him.
\item \textbf{The key is universality} (\S5). Since the value for a single datum can be computed classically, what the
proof needs is the part ``for every initial $\Theta$-datum''.
\item \textbf{The essence is the multiradial algorithm, and the key is its implementation} (\S6). Mochizuki's
description, the results of lana-agents/iut and their reach, and the position of this note.
\item \textbf{The two computations compared as programs, and the multiradial algorithm as a mockup with oracles}
(\S\S7--9).
\item \textbf{Observations, limitations and future work} (\S\S10--12).
\end{enumerate}
Points 1 and 2 are facts that can be checked against earlier results and computations; 3 and 5 are statements of the
sources and our reading of them; 4 is a fact that can be checked by computation together with our reading; 6 is a
proposal of method and a mockup. None of them decides whether IUT is right or wrong.

\subsection{Labels}
\tagS{} a statement of the sources (IUT and Mochizuki's documents); \tagX{} a claim of others (lana-agents/iut, Joshi,
Scholze--Stix, \dots); \tagP{} a fact or computation shown in this note or in the earlier notes; \tagH{} our reading,
interpretation, proposal or position (not a proof).

\section{In every reading, IUT passes through the height inequality}
Corollary~3.12 compares two log-volumes. One is the log-volume $-|\log\qq|$ of the $q$-pilot, which is not subject to
the indeterminacies. The other is the log-volume $-|\log\TT|$ of the holomorphic hull of the union of the possible
images of the $\Theta$-pilot, which is subject to the indeterminacies (Ind1), (Ind2), (Ind3) \cite{IUT3}. The proof of
Theorem~1.10 of IUT~IV bounds $-|\log\TT|$ from above \cite{IUT4}. This gives the chain
\[-|\log\qq|\ \le\ -|\log\TT|\ \le\ U(D).\]
Putting the two ends together, with $a=|\log\qq|$, $k=\ell(\ell+1)/12$ (the average of $j^2$ over the labels
$j=1,\dots,\ell^*$, where $\ell^*=(\ell-1)/2$) and $C=U(D)+k\,a$, gives the height inequality
\[(k-1)\,a\ \le\ C.\]

\tagP{} The value of the middle term $-|\log\TT|$ depends on the ``reading''. The readings have five axes: the range
of the label permutations of (Ind1), container or iteration for (Ind3), the group of (Ind2), whether the hull is taken,
and the archimedean container. \cite{RF} showed the following in each of the 36 readings obtained by combining the axes
found so far.
\begin{itemize}
\item Corollary~3.12 and the upper bound of IV together give the same height inequality $(k-1)a\le C$, because the
two ends $a$, $k$, $C$ are computed by one procedure that takes no reading, and only the middle term depends on the
reading (a conservation law).
\item A reading decides only how much of the wall $(k-1)a$ Corollary~3.12 carries and how much IV carries. No reading
makes both loads nonpositive.
\item In the family obtained by replacing $q$ with $q^N$, the statement that holds for every $N$ carries no positive
load, and the other carries at least the whole wall. Choosing the reading differently does not make the proof cheaper.
\end{itemize}

\tagX{} lana-agents/iut \cite{lana} collects various formulations of the wall and shows that they are equivalent in its
models (\lean{Iut/\allowbreak{}WallEquiv/\allowbreak{}}). For one reading it also shows that Step~(xi) does not follow from the formalized
content of Theorem~3.11 alone (\lean{NonDerivable.lean}; \S6.3). It proves, with no other hypothesis left, that the
variant of Corollary~3.12 for all initial $\Theta$-data gives the abc conjecture (\lean{Iut/\allowbreak{}MainTheorem.lean}).
Earlier, for one reading, Takkun-kohinata/IUT\_LEAN \cite{TK} (15 July 2026) proved in Lean that the containment of
Step~(v) together with Corollary~3.12 gives the height inequality, and conversely (\S1.3).

\tagS{} The route of the sources is Theorem~3.11 $\to$ Corollary~3.12 $\to$ Theorem~1.10 of IV $\to$ ABC. Theorem~1.10
of IV draws its conclusion, in its own statement, by applying the inequality of Corollary~3.12. The sources contain no
route that avoids Corollary~3.12.
\srcquote{\q{and hence, by applying the inequality `$C_\Theta\ge-1$' of [IUTchIII], Corollary 3.12, conclude
that}}{IUT~IV, Theorem~1.10 \cite{IUT4}}
\tagX{} Geiger's audit of this passage describes IV in the same way: there the load of the theta values is carried to the
closed form $(\ell+1)/24$, with that section \q{taking the corollary's conclusion as an input} \cite[\S6.2]{Ge2}
(16 August 2026).

\tagH{} Within the 36 readings of \cite{RF}, whose axes come from the author's notes and in part from lana-agents/iut,
therefore, whoever derives ABC through IUT has to let someone carry the height inequality (its load), in every reading. This is what ``IUT passes through the height
inequality'' means.

\tagX{} There is a position that denies this premise. Joshi holds that the Tate parameter changes together with
(Ind1), (Ind2), (Ind3) \cite{Joshi}. On that view the left end also depends on the reading, and the conservation law
does not apply. This note follows the wording of Corollary~3.12, under which $-|\log\qq|$ is not regarded as subject
to the indeterminacies \cite{IUT3}.

\section{The height inequality is classical}
\tagS{} The inequality that Theorem~1.10 of IV finally gives has the form \cite{IUT4}
\[\tfrac16\cdot\log(\mathfrak q)\ \leq\ \Bigl(1+\tfrac{20\cdot d_{\mathrm{mod}}}{l}\Bigr)\cdot\bigl(\log(\mathfrak
d^{F_{\mathrm{tpd}}})+\log(\mathfrak f^{F_{\mathrm{tpd}}})\bigr)+20\cdot(e_{\mathrm{mod}}^\ast\cdot
l+\eta_{\mathrm{prm}}),\]
where $\log\mathfrak q$ is the normalized degree of the divisor of the $q$-parameters, $\log\mathfrak d$ that of
the different (the normalized logarithm of the discriminant), $\log\mathfrak f$ that of the conductor, and
$d_{\mathrm{mod}}$, $e^*_{\mathrm{mod}}$, $\eta_{\mathrm{prm}}$ are constants.

\tagS{} Mochizuki himself summarizes the result of this computation as a bound on the height of the form expected from
the Szpiro conjecture:
\srcquote{\q{This computation may be summarized as asserting that one obtains a bound on the height of the elliptic
curve under consideration whose leading term is the `log-different $\log(\mathfrak d^K)$' [cf. the final
portion of \S3.7, (iv)] that one expects from [a suitable version of] the Szpiro Conjecture.}}{Mochizuki, The
mathematics of mutually alien copies, \S3.8, $(11^{\mathrm{gau}})$ \cite{Alien}}

\tagP{} In either form, only classical quantities occur in the inequality: the degree of the $q$-parameters, the
discriminant and the conductor, the different of the field $K$, logarithms of primes, and constants. The term $C$ of
the intermediate form $(k-1)a\le C$ of \cite{RF} is also a closed expression in local classical quantities, as observation
(ii) of \S10 shows. To compute both sides for a given curve, one finds $\mathrm{ord}_v(q)$ at each place from the
valuations of the $j$-invariant of the curve and adds up degrees. The program of \S8.2 does just this.

\tagH{} That both sides can be computed is different from the inequality holding for all curves. The latter is a form
of the Szpiro conjecture and is open by methods that do not use IUT. When this note says ``the height inequality is
classical'', it means that \emph{its statement is written in classical quantities only and its value for each datum
can be computed classically}, not that it has been proved classically.

\section{What IUT supplies is a different method of computation}
\tagS{} Mochizuki explains in terms of computation what IUT is a theory of. He casts the problem of heights in the form
``$N\cdot h=h$'' [or the inequality ``$N\cdot h\le h+C$''] and then writes:
\srcquote{\q{One then verifies the desired approximate relation $q^N$ `$\approx$' $q$ by computing
$\{q_{\mathrm{LHS}}^{j^2}\}_j$ --- not in terms of $q_{\mathrm{LHS}}$ [which is immediate from the
definitions!], but rather --- in terms of [the scheme theory surrounding] $q_{\mathrm{RHS}}$ [which is a highly
nontrivial matter!].}}{Mochizuki, The mathematics of mutually alien copies, Introduction \cite{Alien}}
\srcquote{\q{Of course, such an assignment is not compatible with the ring structure of an NF, hence does not exist in
the framework of conventional scheme theory.} \dots{} \q{in some sense the fundamental theme of inter-universal Teichm\"uller
theory consists of the development of a mechanism for computing the effect --- e.g., on heights of elliptic curves
[cf. the discussion of \S2.3!] --- of such non-scheme-theoretic `Gaussian Frobenius morphisms' on NF's.}}{Mochizuki,
The mathematics of mutually alien copies, \S2.4 \cite{Alien}; ``such an assignment'' is
$\qq\mapsto\{\qq^{\,j^2}\}_{j=1,\dots,l^*}$}
\srcquote{\q{Then the $N$-th power map of (i), (ii) corresponds, after passing to heights, to multiplying real numbers
by $N$; the multiradial algorithm corresponds to saying that the height is unaffected (up to a mild error term!) by
multiplication by $N$, hence that the height is bounded!}}{Mochizuki et al., On the formalization of IUT: a preliminary
progress report, \S3 \cite{Form}}

\tagH{} Put side by side, these statements let us read the picture of IUT as follows.
\begin{itemize}
\item If the log-volume of the $\Theta$-pilot $\{\qq^{\,j^2}\}_j$ is computed on the original side
($q_{\mathrm{LHS}}$), it is $-k\,a$ directly from the definitions. This is the ``classical'' computation; it only
multiplies the height by $k$.
\item IUT computes the same quantity on the far side of the $\Theta$-link (the scheme theory surrounding
$q_{\mathrm{RHS}}$), namely in the tensor packets of log-shells, allowing indeterminacies. The result is $-|\log\TT|$,
whose value is $-k\,a$ plus what Mochizuki calls an ``error factor'' (the container, the unit part, the hull and so on:
the remainder $B$ of \cite{RF}).
\item Corollary~3.12 asserts that the value so computed is at least the value $-a$ of the $q$-pilot. Since $k>1$, this
does not hold by the classical computation alone. If it holds, it holds because of the error ($B$) brought in by the
other method of computation. This is the height inequality itself, in the form $N\cdot h\le h+C$.
\end{itemize}

\tagH{} In one sentence, our reading is this. The core of the proof of IUT is to give a ``different way of grasping''
an open inequality (the height inequality) by decomposing both of its sides into a method of computation different
from the existing one, and thereby to attempt a proof that could not be given in the existing framework. In
Mochizuki's own words it is the development of a ``mechanism for computing''. If this reading is right, the content of
IUT lies in the \emph{difference} between the two methods of computation.

\section{The key is universality}
\tagS{} The statement of Corollary~3.12 begins with the sentence
\srcquote{\q{Suppose that we are in the situation of Theorem 3.11.}}{IUT~III, Corollary~3.12 \cite{IUT3}}
The situation of Theorem~3.11 is built from one initial $\Theta$-datum. Corollary~3.12 claims to hold whatever that
initial $\Theta$-datum is.

\tagP{} For one given initial $\Theta$-datum, on the other hand, both sides of the height inequality can be computed
classically (\S\S3, 8). For the six data of this note that are completely factored the inequality holds with a large
margin (Table~2); for D2357 the counted places alone leave it short, and the verdict depends on the places above its
unfactored cofactor \cite[\S5.5]{RF}. What IUT has
to add is therefore not the value for one datum, but the part that it holds for \emph{every} initial $\Theta$-datum.

\tagX{} lana-agents/iut shows in Lean that this ``for every'' is exactly where the strength of ABC lies. The variant of
Corollary~3.12 for all initial $\Theta$-data gives the abc conjecture (\lean{classicalABC\_of\_variant}), and so does
Step~(xi) for all genuine data (\lean{classicalABC\_of\_forall\_xi}).
\srcquote{\q{\textbf{The missing input is ABC-strength}: \lean{classicalABC\_of\_forall\_xi} (Step (xi) for all
genuine data gives \lean{ClassicalABC}, via \lean{classicalABC\_of\_thm311}).}}{lana-agents/iut,
\lean{Iut/\allowbreak{}WallEquiv/\allowbreak{}README.md} \cite{lana}}

\tagX{} Earlier, for one reading, Takkun-kohinata/IUT\_LEAN \cite{TK} (15 July 2026) made the same point: with a ledger
$L=0$ the containment of Step~(v) cannot hold (\lean{no\_lossless\_theta\_hull}), and, since its ledger of IV,
Propositions~1.1--1.4, does not depend on $|\log q|$, supplying that containment uniformly for curves of unbounded height
is equivalent to supplying the height inequality (\S1.3).

\tagP{} The family $q\to q^N$ of \cite{RF} shows the same thing. The statement that holds for every $N$ (the upper bound of
IV) carries no positive load, and the whole wall falls on Corollary~3.12. The same happens in the mockup of \S9 on this
family (\S9.4).

\tagH{} In terms of computer science, Corollary~3.12 is a property over all of infinitely many inputs (a relation
between two outputs). Trying finitely many inputs (testing) never establishes it, however many are tried; an argument
about the program (verification) is needed. IUT seeks to obtain it from the \emph{uniformity} of the construction of
Theorem~3.11, that is, from the fact that the same procedure applies to every initial $\Theta$-datum. The step that does
this is Step~(xi) of the proof of Corollary~3.12. The question ``does IUT prove ABC?'' thus becomes the question
``where does the `for every' of Corollary~3.12 come from?''. The answer of IUT is: from the multiradial algorithm.

\section{The essence is the multiradial algorithm}

\subsection{Mochizuki's description}
\tagS{} The sentence of the progress report quoted in \S4 says that it is the ``multiradial algorithm'' that states
that the height is unaffected by multiplication by $N$ \cite{Form}. The same report also states what the algorithm takes
as input:
\srcquote{\q{Main point: The multiplicative monoid and abstract group structures (but not the ring structures!) are
common (cf. `logical AND $\land$!') to $\dagger$, $\ddagger$ and can be used as the input data for an
algorithm to construct the multiradial representation for the $\Theta$-pilot, i.e., a common container for the
distinct ring strs. (i.e., `arith. hol. strs.')}}{Mochizuki et al., On the formalization of IUT: a preliminary
progress report \cite{Form}}

\tagH{} Restated in terms of computer science, the multiradial algorithm is a procedure that reads only the inputs
common to the two sides (multiplicative monoids and abstract groups) and produces an output (a common container) that
does not depend on the ring structures of the two sides. If the ring structures are regarded as secret inputs, this is
the requirement of \emph{non-interference} in information-flow security (the output does not depend on the secret
inputs). The mockup of \S9 expresses this separation in the types of the program.

\tagS{} Mochizuki himself uses the phrase ``mutually non-interference'' in a mathematical sense, for the Kummer
theories on the two sides of the $\Theta$-link:
\srcquote{\q{the mutually non-interference or simultaneous executatibility of the Kummer theories surrounding the
$q$-pilot and $\Theta$-pilot} \dots{} \q{when the Kummer theory surrounding the $q$-pilot is held fixed, and one allows
the Kummer theory surrounding the $\Theta$-pilot to be subject to various indeterminacies.}}{Mochizuki, On the
essential logical structure of IUT, \S3.5, (NonInf) \cite{EssLgc}}

\tagX{} The interim report of Project LANA lists as open the question whether the value-group part affects the output
of the algorithm:
\srcquote{\q{As suggested here, the input to the `raw' Theorem 3.11 algorithm appears to consist solely of the
\'etale-unit BPS, and the value-group part does not seem to affect the algorithm's result.} \dots{} \q{Reconciling
these points remains a topic of discussion among LANA members.}}{LANA interim report, \S6.2, Remark~6.2.2
\cite{LANA}}

\tagH{} The ``non-interference'' of this note is the property of information-flow security (the output does not depend
on secret inputs). Mochizuki's (NonInf) states that two Kummer theories are compatible; we do not claim that the two
are the same. The open question of LANA, on the other hand, has exactly the form of a question of non-interference:
whether the output depends on an input, the value group. The mockup of \S9 puts questions of this kind into a form
that can be posed in types and checks.

\subsection{The position of this note}
\tagH{} This note regards the multiradial algorithm as the essence of IUT, for four reasons.
\begin{enumerate}
\item It is this algorithm that actually performs the ``different method of computation'' of \S4: its content is to
compute the $\Theta$-pilot again on a side that does not share the ring structure.
\item The universality of \S5 is supposed to come from the uniformity of this algorithm. Corollary~3.12 begins from
``the situation of Theorem 3.11'', and Step~(xi) makes the comparison.
\item The other steps are classical or elementary. The estimate of IV is a local computation of log-shells and
differents (\S8.3), and the route from Corollary~3.12 to ABC has been formalized in Lean by lana-agents/iut (\S2).
\item The C program of IUT in \S8 only computes the estimate of IV of the output of this algorithm; it is not the
algorithm. Both sides come out without it. The part that could not be written in the program is exactly the part of
the algorithm.
\end{enumerate}

\tagS{} Mochizuki regards the step from Theorem~3.11 to Corollary~3.12 (``Stage~1'') as the essential nontrivial content
of the theory. He also writes that one can understand the delicate inter-relations of the parts of the structures only
by working with specific algorithms:
\srcquote{\q{We chose to concentrate on this aspect of the theory first since this aspect of the theory --- i.e.,
`Stage 1: [IUTchIII] Theorem 3.11 $\Longrightarrow$ Corollary 3.12' (cf. the discussion at the beginning of \S2) ---
has received the most public attention. Indeed, this aspect corresponds to the essential nontrivial content of the
theory, i.e., that the height of the elliptic curve under consideration is equal to $N$ times the height of the elliptic
curve (where $N$ is a large positive number), up to a small discrepancy (arising from (Ind1,2,3) + hull), thus implying a
bound on the height.}}{Mochizuki et al., On the formalization of IUT: a preliminary progress report \cite{Form}}
\srcquote{\q{That is to say, although different algorithms, or software, may yield the same output data, it is only by
working with specific algorithms that one may understand the delicate inter-relations that exist between various
components of the structures that occur as one performs various operations [i.e., the mutations of a
mutation-history].} \q{In the case of the theory developed in the present series of papers, one central example of
this phenomenon is the cyclotomic rigidity isomorphisms that underlie the theory of
$\Theta^{\times\boldsymbol\mu}_{\mathrm{LGP}}$-link compatibility discussed in [IUTchIII], Theorem 3.11, (iii),
(c), (d)}}{IUT~IV, Remark~3.6.2, (i) \cite{IUT4}}

\tagH{} Mochizuki takes Stage~1 as the essential content, and this note takes the multiradial algorithm, which carries
it, as the essence. The difference is one of emphasis; we agree that what should make Stage~1 hold is the uniformity of
the algorithm (\S5). Mochizuki also names the cyclotomic rigidity isomorphisms as the central example of the relations
that only specific algorithms reveal. The first thing implemented in the mockup of \S9 is a finite model of cyclotomic
rigidity.

\tagH{} This note therefore takes the following position. \textbf{One has not understood IUT essentially unless the
multiradial algorithm has been elucidated and implemented.} Here an ``implementation'' means giving the procedure that
computes the output from the data common to the two sides alone, either in a form that runs on finite data or as a
construction that can be verified formally. Either way, whether the output satisfies Step~(xi) can then be asked of the
implementation.

\subsection{The results of lana-agents/iut and their reach}
\tagX{} lana-agents/iut has shown that Step~(xi) does not follow from the formalized content of Theorem~3.11:
\srcquote{\q{\textbf{Theorem 3.11 of IUT III, as formalised here, does not imply Step (xi) of the proof of
Corollary 3.12.}} \dots{}
\q{a proof of Step (xi) for all genuine \lean{D} must use facts about \lean{D} beyond them, and such a proof would give
the ABC conjecture}}{lana-agents/iut, the module docstring of \lean{Iut/\allowbreak{}WallEquiv/\allowbreak{}NonDerivable.lean} \cite{lana}}

\tagX{} Over its package, Step~(xi) is equivalent to the Scholze--Stix inequality
$(s-1)\,|\log q|\le\text{log-vol}(\text{hull}(VC(I)))$ (\lean{xi\_iff}; $s$ is the $k$ of this note). This is itself a
height inequality, with the volume of the container on the right. Moreover the package is realized for all values of
$a$ and $V$ (\lean{family}). Hence no inequality bounding $a$ follows from the formalized content of Theorem~3.11 (that
the package exhausts what is proved at the level of Theorem~3.11 is, as lana-agents/iut itself says, audited and not a
theorem). lana-agents/iut also writes:
\srcquote{\q{The repository proves no other relation coupling \lean{V} with \lean{a}.} \dots{}
\q{\textbf{Not shown}: that Step (xi) fails for some genuine \lean{D}, or that ABC fails; the package is a
model}}{lana-agents/iut, the module docstring of \lean{Iut/\allowbreak{}WallEquiv/\allowbreak{}NonDerivable.lean} \cite{lana}}

\tagX{} The plan for formalizing the whole of Theorem~3.11 in the same repository (a draft of 14 October 2026) says that
what exists is a theorem about a model and not the algorithm, and that a genuine Theorem~3.11 would not by itself decide
Step~(xi):
\srcquote{\q{We have IUT III Thm 3.11 as a theorem \emph{about a model}. We do not have the algorithm.} \dots{}
\q{It does not decide Cor 3.12, Step (xi).} \dots{}
\q{A genuine 3.11 would turn that from a statement about our models into a statement about IUT's actual objects. It
would not prove or refute SS.}}{lana-agents/iut, \lean{Iut/\allowbreak{}Multiradial/\allowbreak{}ROADMAP.md}, \S0 \cite{lana}}

\tagH{} From the above, this note reads the non-derivability of lana-agents/iut as follows.
\begin{itemize}
\item What lana-agents/iut has shown is \emph{the result of taking the content of Theorem~3.11, as currently written,
as a specification}. Neither Step~(xi) nor any height inequality follows from that specification. But this is not a
proof that IUT has suffered fatal damage that cannot be repaired. lana-agents/iut itself writes that it has not shown
that Step~(xi) fails for some genuine datum.
\item This note holds that the key lies in the implementation of the multiradial algorithm. Step~(xi) can be asked of
the actual objects of IUT, rather than of a model, only once an implementation exists (the ``statement about IUT's
actual objects'' of the plan of lana-agents/iut).
\item The strength of what would have to be found there is, however, fixed: a uniform property of ABC strength that is
not written in the specification. In the words of lana-agents/iut, a relation coupling the size $V$ of the possible
images with $a=|\log\qq|$, or facts about genuine data beyond the package. One may not expect an implementation to fill
a small gap; an implementation only makes visible whether the actual construction has such a property. In the mockup
of \S9, the place where it would enter is the oracle \orc{step-xi}.
\item There is an opposite position. Scholze--Stix hold that if the identifications of the copies of the real numbers
are taken consistently, the inequality of Corollary~3.12 becomes empty \cite[\S2.2]{SS}. On that view, an
implementation will not produce the required property.
\end{itemize}

\tagX{} Earlier, Higuchi \cite{Hi} (31 March 2026) and Pajo \cite{Pj} (26 April 2026) proved that the inclusion of
Step~(xi-f), or the numerical inequality drawn from it, does not follow from qualitative or scale-free abstractions of
this passage. \tagP{} The author's note \cite{W311} proves a
statement of the same kind from the results that the proof of Corollary~3.12 cites: their log-volume content, stated as
axioms at the strength of the source, has models in which the conclusion fails, one of them on actual local objects
\cite[Theorems~10 and~22]{W311}. By \cite[Proposition~6]{W311}, the carrier reading of (SHE) (the $q$-pilot is one of
the possible images) implies the conclusion, and the value reading (its log-volume is one of the possible values) is
equivalent to it. \tagH{} In the mockup of \S9 these are what the oracle \orc{step-xi} returns; an implementation would
have to supply one of them.

\section{Method: comparing the two computations as programs}
\tagH{} To study the difference between the two methods of computation, we propose to place them side by side as two
programs that read the same input, instead of comparing them in the language of difficult mathematics. There are three
reasons.
\begin{enumerate}
\item \textbf{The flow of data becomes explicit.} A program writes, line by line, what it reads, what it builds, and
what it compares where. A reader who does not share the conventions of the mathematics can follow which input affects
which output.
\item \textbf{The two are compared on the same ground.} If the two programs are given the same input (the local data of
the same curve), the only difference is the path of the computation.
\item \textbf{The tools of computer science can be used}, for example:
\begin{itemize}
\item program equivalence and loop summarization (symbolic evaluation): are the two computations the same function of
the input, and if not, where do they differ?
\item analysis of data flow and information flow: along which paths does the size of $q$ reach the output?
Multiradiality can be written as non-interference with respect to a secret input, the ring structure (\S6.1).
\item relational verification: a framework that treats claims of the form ``for all inputs, the output of program A
$\le$ the output of program B'' as something to prove rather than to test. Corollary~3.12 has this form (\S5, \S10
(v)).
\item mockups with oracles: the parts that cannot be implemented are named as oracles, and where their values come from
is recorded, so that the code distinguishes what is implemented from what is not (\S9).
\item adversarial families of inputs: instead of checking on inputs that happen to have a large margin, one tries
families of inputs that make the wall large (the family $q\to q^N$ of \cite{RF}).
\end{itemize}
\end{enumerate}

\tagS{} Mochizuki himself has a caveat on what this method cannot reach. When the Gaussian integral is computed in
polar coordinates, it is a hopeless task to identify which portion of the original one-dimensional integral
``corresponds'' to the radial integral. In the same way he writes:
\srcquote{\q{it is essentially a hopeless task to trace, in some sort of explicit or readily computable fashion, the way
in which the value group portions $(b^\Theta)$, $(c^\Theta)$ in the domain of the $\Theta$-link [cf.
$(9^{\mathrm{gau}})$] appear within the multiradial representation via tensor packets of log-shells [cf.
$(10^{\mathrm{gau}})$] expressed in objects arising from the codomain of the $\Theta$-link}}{Mochizuki, The
mathematics of mutually alien copies, \S3.8 \cite{Alien}}

\tagH{} This is not a statement that the tools of computation are insufficient, but a structural statement: the
corresponding pieces do not exist in any meaningful form. This note therefore does not look for a line-by-line
correspondence between ``this line'' of the classical computation and ``this line'' of the computation of IUT. What we
compare is the shape of the whole computation: what is read as input, what intermediate structures are built, where $q$
enters, and which comparison remains, uncomputed, as a claim. What Mochizuki says cannot be traced is the content of the
comparison of Corollary~3.12. In the mockup of \S9 it lies inside the oracle \orc{step-xi}.

\section{The two programs}

\subsection{The common input}
\tagP{} The input is the local data of one initial $\Theta$-datum: the prime $\ell$, the field $F=\Q(\sqrt D)$, and, for
each place of $F$ above a prime dividing $N(a)N(1-a)$ (for the curve $y^2=x(x-1)(x-a)$), the prime $p$, the residue
degree $f$, the ramification index $e_F$ of $F_v/\Q_p$ and $\mathrm{ord}_v(q)$ (0 for good reduction). For example, the
file for D211 begins
\datafile
The local data were taken from an elementary program that does not use IUT (\lean{elementary\_height.py}, which finds
$\mathrm{ord}_v(q)$ from the valuations of the $j$-invariant). We checked that the places it gives agree, prime by
prime, as multisets of $(f,e_F,\mathrm{ord}\,q)$, with those built independently by the program of the walkthrough \cite{WT},
for all seven data (\lean{export\_local\_data.py}). The two C programs and the C++ mockup all read this same file.

\subsection{The classical computation}
\tagP{} The classical computation uses only the places of bad reduction, their residue degrees and $\mathrm{ord}_v(q)$.
No labels, tensor packets or log-shells occur. The constants are copied from the statement of Theorem~1.10; its proof
is not used.
\VerbatimInput{../c/classical.c}
To compare with the IUT side, the places counted are those of $\Vbad$ of the catalogue \cite{Cat} minus those above 3 and 5
(as in \cite[\S5]{Gap}, and \cite{WT}; Definition~3.1, (b), of IUT~I allows a subset \cite{IUT1}). The values over all places of bad
reduction are printed too.

\subsection{The computation of IUT}
\tagP{} From the same input, the computation of IUT first builds the structure that IUT attaches to each place:
\begin{itemize}
\item the weight $w_v=[F_v:\Q_p]/[F:\Q]$ (summing to 1 over the places above $p$);
\item membership in $\Vbad$;
\item the ramification index $e_K=e_F\cdot15\ell/\gcd(15\ell,\mathrm{ord}\,q)$ of $K=F(E[\ell])$, with
$F=F_{\mathrm{mod}}(\sqrt{-1},E[15])$ (at bad places; above $p$ the formula of Appendix~A of the catalogue, assumed at
the other bad places);
\item the order of $\qq=q^{1/2\ell}$, the outer radius $\rho$ of the log-shell $(1/2p)\log(\Ocal_v^\times)$, and the
order $(e-1)/e$ of the different.
\end{itemize}
Then, for each prime carrying theta values, it enumerates the labels $j=1,\dots,\ell^*$ and the tuples
$t=(t_0,\dots,t_j)$ of places, one for each factor of the tensor packet (weight $w(t_0)\cdots w(t_j)$), and for each
tuple computes
\begin{itemize}
\item the hull on the $\Theta$-side (the reading of IV: the theta value sits on the factor $j$; container, units,
hull, small archimedean part): $\sum_i\rho(t_i)-j^2\,\mathrm{ord}(\qq_{t_j})$;
\item the local estimate of Proposition~1.4, (iii), of IV with $\lambda=j^2\,\mathrm{ord}(\qq_{t_j})$:
$(-\lambda+\sum_i\mathfrak d(t_i)+1)\log p+\sum_{e(t_i)>p-2}(3+\log e(t_i))$.
\end{itemize}
These are averaged with the weights of the tuples and over the labels (the procession normalization). The places
above $2$, $3$, $5$, $\ell$ are not computed: on the side of the hull their contribution, which is nonnegative, is set
to 0, and on the side of the upper bound they are bounded by Step~(ii) (the different) and (R4) (the terms of $I^*$).
Finally the archimedean part $\frac{\ell+5}4\log\pi$ is added, and the program goes on to Step~(viii) and the form of
Theorem~1.10. The central enumeration:
\iutcore

\subsection{Output and cross-check}
\begin{center}\small
\begin{tabular}{@{}lrrrrrrr@{}}
\toprule
& \multicolumn{2}{c}{classical} & \multicolumn{5}{c@{}}{IUT}\\
\cmidrule(lr){2-3}\cmidrule(l){4-8}
datum & $\log\mathfrak q$ & $\log\mathfrak f$ & $-|\log\qq|$ & $-|\log\TT|$ (IV) & $U(D)$ & wall $(k-1)a$ & $C$\\
\midrule
\tabletwo
\bottomrule
\end{tabular}

\smallskip Table 2: output of the two programs (counted places). The values for D2357$^\ast$ are conditional, since
the places above its unfactored cofactor are not counted; for its counted places the wall exceeds $C$. That D2357 is an
initial $\Theta$-datum satisfying the hypotheses of Theorem~1.10 of IUT~IV is proved in \cite[Theorem~2.1]{Fourteen}.
\end{center}
\tagP{} For every datum, $2\ell\cdot|\log\qq|$ on the IUT side agrees with $\log\mathfrak q$ on the classical side, and
the two sides of the inequality of Theorem~1.10 that the two programs finally print agree (the right-hand side is
about $1.548\times10^8$ for $\ell=7$ and about $3.47\times10^9$ for D2357). We compared the output of the two C programs
with the Python programs of the bundle of \cite{RF} (\lean{elementary\_height.py}, \lean{iut\_height.py},
\lean{iut\_readings36.py}) on the seven data: 119 items agree within a relative error of $10^{-9}$
(\lean{check\_c.py}). The items are, on the classical side, $\log\mathfrak q$ and $\log\mathfrak f$ (over all places
and over the counted ones), $\log\mathfrak d$ and the right-hand side of the theorem; on the IUT side $-|\log\qq|$,
$-|\log\TT|$, $U(D)$ with its parts, Step~(viii) and the right-hand side of the theorem; and the two ends against each
other.

\section{A mockup of the multiradial algorithm with oracles}
\tagP{} The program of IUT in \S8 only computes the estimate of IV of the output of the multiradial algorithm. The C++
mockup of this section makes a place for the algorithm. The parts that can be implemented at present are implemented,
and the rest is left as oracles. The values of the oracles are reverse-computed from the ring structure of a theater or
from the conclusion. That they are reverse-computed is written in the comments of the code, in the marks of the output
(such as \texttt{[ORACLE step-xi]}) and the word \texttt{REVERSE-COMPUTED}, and in a ledger at the end of the output. The
oracles are named by their content rather than numbered (Table~5). \textbf{The output of the mockup proves nothing.}

\subsection{Design}
\tagP{} The mockup builds two Hodge theaters $\dagger$, $\ddagger$ from the same initial $\Theta$-datum. The ring side of
a theater (the ramification of $K_v$, the valuations of $q$, and the coordinates that the theater chose on its cyclotome
$\Delta_\Theta$) is private. The implemented part of the multiradial algorithm receives only the data common to the two
sides and the \'etale-like data of $\dagger$. The header \texttt{common.hpp} that defines them does not know the class of
the theater. The file of the algorithm, \texttt{multiradial.cpp}, does not include the header of the theater, even
indirectly, so code that reads the ring of a theater there does not compile.
\begin{center}\small
\begin{tabularx}{\textwidth}{@{}l X >{\raggedright\arraybackslash}p{3.6cm}@{}}
\toprule
file & role & reads the ring of a theater?\\
\midrule
\texttt{common.hpp} & everything that crosses from a theater to the algorithm: types of abstract groups, Kummer classes,
boxes of theta groups, global data & no\\
\texttt{theater.*} & the ring side of a Hodge theater; all private & itself\\
\texttt{oracle.*} & the six oracles; they reverse-compute their values and enter them in a ledger & yes (friends of the
theater)\\
\texttt{multiradial.*} & the implemented part (A)--(G) of the multiradial algorithm & cannot (headers checked
transitively)\\
\texttt{iv.*} & the estimate $U(D)$ of IV, Step~(viii), the form of Theorem~1.10 & cannot\\
\texttt{link.*} & the $\Theta$-link: checks that the data common to the two theaters agree place by place, and links
the $\Theta$-pilot of $\dagger$ to the $q$-pilot of $\ddagger$; not a map of rings & only through the oracles\\
\texttt{classical.hpp} & the classical computation (as in \S8.2) & reads the local data of the curve directly\\
\texttt{main.cpp} & report, check of non-interference, scan of the family $q^N$ & ---\\
\bottomrule
\end{tabularx}

\smallskip Table 3: the files of the mockup (\cppfiles{} files, \cpplines{} lines in all).
\end{center}
The flow of the computation is: initial $\Theta$-datum $\to$ theaters $\dagger$, $\ddagger$ (each choosing its own
coordinates) $\to$ $\Theta$-link $\to$ oracles \orc{gab}, \orc{global}, \orc{kummer}, \orc{theta-group}
(reverse-computing the \'etale-like data of $\dagger$) $\to$ implemented part (A)--(G) $\to$ oracle \orc{ind3} $\to$
the $q$-pilot of $\ddagger$ $\to$ oracle \orc{step-xi} (reverse-computing Step~(xi)) $\to$ Corollary~3.12 $\to$ the
estimate of IV $\to$ the height inequality and Theorem~1.10. What reaches the algorithm (a part of
\texttt{common.hpp}):
\cppbox

\subsection{What is implemented}
\begin{center}\small
\begin{tabularx}{\textwidth}{@{}l X >{\raggedright\arraybackslash}p{4.2cm}@{}}
\toprule
step & what it does & basis\\
\midrule
(A) & recovers $p,d,f,e$ from the type of the abstract group
$G_{K_v}^{\mathrm{ab}}\cong\widehat\Z\times\Z_p^d\times(\text{torsion})$: $p$ is the prime with an extra $\Z_p$-part,
$d$ its rank, $f$ is read off from the order $p^f-1$ of the torsion prime to $p$, and $e=d/f$ & local class field
theory; the plan of lana-agents/iut also lists the local invariants of $G_K$ as already genuinely constructed\\
(B) & the outer radius $\rho$ of the log-shell and the order of the different, from $(p,e)$ & the log-shells of IV,
\S1 (the formulas of \S8.3)\\
(C) & cyclotomic rigidity at level $N=2^{61}-1$: recovers the coordinates that the theater chose on $\Delta_\Theta$
from the commutator $[t,a]$ in the box of the theta group & the remark of the paper on the \'etale theta function
quoted below, and the Heisenberg group of \lean{Iut/\allowbreak{}EtaleTheta/\allowbreak{}Heisenberg.lean} of lana-agents/iut\\
(D) & reads the valuations of the theta values $\qq^{\,j^2}$ from their Kummer classes written in the coordinates of
$\Delta_\Theta$, and checks the structure $\kappa(\theta_j)=j^2\kappa(\theta_1)$ & Kummer theory at a finite level (the
units are $N$-th powers)\\
(E) & the procession: labels $j=1,\dots,\ell^*$ and the tuples of places of the tensor packets & as in \S8.3\\
(F) & possible images and their hull; in the reading of IV, (Ind1) $=$ average over the labels, (Ind2) $=$ units,
(Ind3) $=$ container & as in \S8.3\\
(G) & the archimedean container $\frac{\ell+5}4\log\pi$ & as in \S8.3\\
\bottomrule
\end{tabularx}

\smallskip Table 4: the implemented part (\texttt{multiradial.cpp}).
\end{center}
\tagS{} Step (C) is a model, at a finite level, of the following point that the paper on the \'etale theta function
makes about cyclotomic rigidity:
\srcquote{\q{One further interesting aspect of cyclotomic rigidity in the theory of the present paper is that it is
obtained essentially as the result of a certain computation involving commutators [cf. Remark 2.19.2]. Put another
way, one may think of this cyclotomic rigidity as a sort of consequence of the nonabelian structure of [what
essentially amounts, from a more classical point of view, to] the `theta-group'.}}{Mochizuki, The \'etale theta
function and its Frobenioid-theoretic manifestations, Introduction \cite{EtTh}}
\tagP{} The theta group of the mockup is the product $(c,a,b)(c',a',b')=(c+c'+b\,a',\,a+a',\,b+b')$, written in the
coordinates that the theater chose on $\Delta_\Theta$. \tagX{} This is the Heisenberg group of lana-agents/iut
(\lean{ThetaGroup} in \lean{Iut/\allowbreak{}EtaleTheta/\allowbreak{}Heisenberg.lean}, added on 12 October 2026), whose
multiplication is $(c+c'+(l\cdot b)\,a',\,a+a',\,b+b')$, in the case $l=1$ \cite{lana}. The algorithm may use
only multiplication, inversion and the two generators. The $c$-coordinate of the commutator $[t,a_1]$ gives the unit
$\gamma$ of the theater's coordinates directly.
\cpprigidity

\subsection{What is left as oracles}
\begin{center}\small
\begin{tabularx}{\textwidth}{@{}>{\raggedright\arraybackslash}p{2.45cm} *{3}{>{\raggedright\arraybackslash}X}@{}}
\toprule
name & content in IUT & how the mockup fills the value & why it cannot be implemented now\\
\midrule
\orc{gab} & the abstract group $G_v$ (an infinite profinite group) & reverse-computes the type
$(p,[K_v:\Q_p],p^f-1)$ from the field $K_v$ by local class field theory & there is no way to hand over a profinite group
as finite data; the mockup hands over the type itself, so $p$ is written in the input\\
\orc{global} & the copies of $F_{\mathrm{mod}}$ (global Frobenioids) and their Kummer theory: weights, labels of
primes, $\log\mathfrak d^F$, $\log\mathfrak f^F$ & reverse-computed from the number field $F$ and $\Vbad$ & depends on
global anabelian reconstruction (Belyi cuspidalization); by the plan of lana-agents/iut it presupposes the relative
$p$-adic Grothendieck conjecture, whose proof uses $p$-adic Hodge theory\\
\orc{kummer} & the Kummer class of the \'etale theta function and the Kummer classes of the theta values obtained by
its Galois evaluation; the Kummer classes of the $q$-pilot of $\ddagger$ & reverse-computed at level $N$ from the
valuations of $q$ in the field (the theta values written in the coordinates of $\Delta_\Theta$) & tempered fundamental
groups and the \'etale theta function cannot be held as finite objects\\
\orc{theta-group} & mono-theta environments (the theta quotient of the tempered fundamental group) & builds the box of
a Heisenberg group at level $N$ from the theater's coordinates on $\Delta_\Theta$ & the same; but the computation of the
rigidity from the box, step (C), is implemented\\
\orc{ind3} & upper semi-compatibility of (Ind3) (IUT~III, Proposition~3.5, Theorem~3.11, (ii), (a)) & only assumed
(the possible images are the containers) & no iteration of log-links and no log-Kummer correspondence\\
\orc{step-xi} & Step~(xi) of the proof of Corollary~3.12: the $q$-pilot of $\ddagger$ lies in the hull of the possible
images & returns the result of comparing the two log-volumes of the datum: \textbf{reverse-computed from the
conclusion} & this is the disputed part; it does not follow from the formalized content of Theorem~3.11 (\S6.3)\\
\bottomrule
\end{tabularx}

\smallskip Table 5: the oracles (\texttt{oracle.cpp}; every value is reverse-computed).
\end{center}
\tagP{} The oracle for Step~(xi) is shown below. What it compares is the conclusion itself, and the comment and the
returned value say so.
\cppstepxi

\subsection{Checks and output}
\begin{itemize}
\item \tagP{} \textbf{Cross-check.} On the seven data the output of the mockup was compared with that of the C programs;
\cppitems{} items agree within a relative error of $10^{-9}$ (\lean{check\_cpp.py}). Since the C programs agree with
the Python programs on 119 items (\S8.4), the mockup agrees with the Python programs too. The items are the two ends,
$-|\log\TT|$, $U(D)$ with its parts, the wall, $C$, Step~(viii), the two sides of Theorem~1.10, the values of the
classical side, and, prime by prime, the enumerated and the closed values of the hull.
\item \tagP{} \textbf{Information flow.} Following the headers transitively with \texttt{g++ -H}, we checked that
\texttt{multiradial.cpp} and \texttt{iv.cpp} include the headers of the theater, the oracles and the link not even
indirectly (\texttt{make check}).
\item \tagP{} \textbf{Non-interference.} When the coordinates that the theaters choose on $\Delta_\Theta$ (a unit of
$\Z/N$ per place) are changed in eight ways, the output is the same bit for bit (all seven data). Without step (C) the
output depends on that choice ($-|\log\TT|\approx\selftestnaive$ for D211, and so on), although the $j^2$-structure is
still seen.
\item \tagP{} \textbf{The family $q\to q^N$.} We ran $N\le40$ prime to $\ell$ (Table~6). The implemented part runs the
same steps for every $N$. The reverse-computed Step~(xi) (the oracle \orc{step-xi}) answers ``yes'' only for small $N$,
for all six data; the height inequality holds up to much larger $N$. This family consists in general not of initial
$\Theta$-data of actual curves (it is an adversarial input).
\end{itemize}
\begin{center}\small
\begin{tabular}{@{}lll@{}}
\toprule
datum & $N$ for which \orc{step-xi} answers ``yes'' & $N$ for which the height inequality holds\\
\midrule
\tablesix
\bottomrule
\end{tabular}

\smallskip Table 6: the family $q\to q^N$ ($N\le40$, $N$ prime to $\ell$). D2357 was not run; for it the oracle
\orc{step-xi} answers ``no'' already at $N=1$ (Table~2).
\end{center}
The output of the mockup for D211, its check of non-interference and its scan of $q\to q^N$ for $N\le20$ are in
Appendix~A.

\section{Observations}
What follows is what one sees by reading the two programs and the mockup side by side. None of it decides whether IUT is
right or wrong.

\paragraph{(i) Input, the structures built, and the two ends.} \tagP{} The two programs start from the same input. The
IUT side builds many structures from it: the ramification indices of $K$, weights, labels, tuples of places, radii of
log-shells, differents, and the estimates above $2,3,5,\ell$. But all of them are determined by the input, and no new
input is read. The numbers at the two ends ($\log\mathfrak q$ and the right-hand side of Theorem~1.10) agree (\S8.4). The
classical side reads only $f$ and $\mathrm{ord}\,q$ at the places of bad reduction and the discriminant of $F$; the
places of good reduction and $e_F$ are read only by the IUT side (to build the weights and the ramification of $K$).

\paragraph{(ii) Folding the right-hand side.} \tagS{} In the reading of IV, every term of the enumeration of tuples is a
sum of terms that depend on one factor only, and the weighted average folds into a closed expression. This folding is
printed in the proof of Theorem~1.10 of IUT~IV itself, Step~(v): Proposition~1.7 of IV gives the weighted average over
the tuples, and the identities (E1), (E2) give the average over the labels, with the coefficients $\frac{\ell+5}4$ on
$\log\mathfrak d^K$ and $\frac{\ell+1}{24}$ on $\log\mathfrak q$.
\srcquote{\q{Thus, by applying Proposition 1.7, we obtain that the resulting `weighted average upper bound' is given
by} \dots{} \q{By averaging over $j\in\{1,\dots,l^{\divideontimes}=\frac{l-1}2\}$ and applying (E1), (E2), we obtain
the `procession-normalized upper bound'\,}}{IUT~IV, proof of Theorem~1.10, Step~(v), pp.~28--29 \cite{IUT4}}
Remark~1.10.1 of IV observes that this computation of the leading term via (E1), (E2) is essentially identical to the
computation of the leading term in the proof of Theorem~A of Mochizuki's earlier Hodge--Arakelov theory, and that
\q{the computations performed in the proof of Theorem 1.10 were already essentially known to the author around the year
2000; the problem then was to construct an appropriate framework, or geometry, in which these computations could be
performed!} \cite[Remark~1.10.1]{IUT4}.

\tagP{} In the notation of the programs, for each prime $p$,
\begin{align*}
U_p&=\tfrac{\ell+5}4\Bigl(\log\mathfrak d^K_p+\sum_{e_v>p-2}w_v(3+\log e_v)\Bigr)-\tfrac{\ell+1}{24}\log\mathfrak
q_p+\log p,\\
M_p&=\tfrac{\ell+5}4\sum_vw_v\rho_v\log p-\tfrac{\ell+1}{24}\log\mathfrak q_p,
\end{align*}
where $\tfrac{\ell+5}4$ is the average over the labels of the number $j+1$ of factors and $\tfrac{\ell+1}{24}\log\mathfrak
q_p=k\,a_p$. The C program prints $U_p$, and the mockup prints both $U_p$ and $M_p$, as enumerated and in closed form;
the two agree on the six data with $\ell=7$ (for D2357 the long tuples are counted by linearity). \tagX{} Dupuy--Hilado
carry out the same averaging as an iterated expectation over tuples and labels, with the average $\frac{\ell+5}4$ of the
number of factors \cite{DH}.

\tagH{} That is, as a function of the input, the right-hand side of the computation of IUT is a closed expression in
local classical quantities ($\log p$, ramification indices, $\mathrm{ord}\,q$, weights). The enumeration of labels and
tuples makes no new numbers; it rearranges classical quantities. The intermediate form $(k-1)a\le C$ is, in this sense,
an inequality in classical quantities only. This holds in the reading of IV. In other readings (L, L0, which enlarge the
label permutations), terms taking a minimum enter, and the closed expression becomes an order statistic of the
distribution of the orders of $q$ (\cite[Proposition~13.2]{II}; \cite{RF}).

\paragraph{(iii) How $q$ enters, and the coefficient of the wall.} \tagP{} On the classical side, $q$ enters as
$\log\mathfrak q$ directly. On the IUT side, the size of $q$ enters in two places: the $q$-pilot
$-|\log\qq|=-\log\mathfrak q/(2\ell)$, and the theta values $j^2\,\mathrm{ord}(\qq)$. Averaged over the labels, the
latter becomes $-k\,a$, and enters the hull of the $\Theta$-side and the upper bound of IV in the same form. The
remaining terms (log-shells, differents, $I^*$, $2,3,5,\ell$, the archimedean part) depend on the orders of $q$ only
through where the reduction is bad, whether $\ell$ divides them, and $\gcd(15\ell,\mathrm{ord}\,q)$ (the ramification of
$K$), not through their size.

\tagH{} Hence the coefficient of the wall $(k-1)a$ is decided by the average $k$ of $j^2$ over the labels alone. The
classical side has nothing corresponding to this $k$. The ``Gaussian Frobenius morphism'' $\qq\mapsto\{\qq^{\,j^2}\}_j$
of Mochizuki (an assignment not compatible with the ring structure) appears in the program as this one coefficient $k$.

\paragraph{(iv) Where the constant 6 comes from.} \tagP{} The constant $\frac16$ on the classical side is copied from
the statement of Theorem~1.10. On the IUT side, the ratio of the coefficient $\frac{\ell+5}4$ of the leading term of
the margin (on $\log\mathfrak d^K$) to the coefficient $(k-1)/(2\ell)$ of the wall (on $\log\mathfrak q$) gives
\[c_\ell=\frac{\ell+5}4\cdot\frac{2\ell}{k-1}=\frac{6\ell(\ell+5)}{(\ell+4)(\ell-3)},\]
in the form $\log\mathfrak q\le c_\ell\log\mathfrak d^K+\cdots$, with $c_7=11.45$, $c_{157}=6.15$, and $6$ in the limit
$\ell\to\infty$. \tagX{} This constant is in Dupuy--Hilado, as $\frac1{6+\varepsilon_\ell}$ with
$\varepsilon_\ell=(24\ell+72)/(\ell^2+\ell-12)$, where they also note $\varepsilon_\ell=O(1/\ell)$ \cite{DH};
$6+\varepsilon_\ell=c_\ell$. The estimate of Zhou has the same factor $(\ell^2+5\ell)/(\ell^2+\ell-12)=c_\ell/6$
\cite{Zhou}, and the Lean port of motizukipdf proves the bound $\log\mathfrak q\le c_\ell\log\mathfrak d^K$ from the
comparison in the common container (\lean{corThreeTwelve\_unfolded}, with $c_\ell$ written as
$\frac{12\ell(\ell^*+3)}{\ell(\ell+1)-12}$) \cite{Richards}. Observation (iv) is not a new result; it rereads this in the
language of programs. \tagP{} The two ingredients are computed in the author's earlier reports: the average of the
stage coefficients is $\ell(\ell+1)/12$ times the stage-one coefficient \cite[Proposition~10]{Fourth}, and
$12/(\ell(\ell+1))$ is compared with the $12/\ell^2$ used in the rearrangement in Step~(viii) of the proof of
Theorem~1.10 \cite[Propositions~1--2]{Fifth}.

\tagH{} In terms of programs, the constant 6 is the ratio of the averages of two loops: the average over the labels $j$
of the number $j+1$ of factors (how many times the term of the different is added), and the average of the weights
$j^2$ of the theta values minus 1 (with what multiple $q$ enters).

\paragraph{(v) The comparison that remains uncomputed.} \tagP{} The two programs only compute the two sides for a given
datum; they do not prove the inequality. In the programs, Corollary~3.12, the core of the proof of IUT, is just the
one-line comparison $-|\log\qq|\le-|\log\TT|$. In the reading of IV, the value $M=-|\log\TT|$ on the $\Theta$-side
(the sum over the primes of the $M_p$ of (ii), plus the archimedean part $\frac{\ell+5}4\log\pi$ and the contribution
above $2,3,5,\ell$) is $-k\,a+B$ with $B=M+k\,a=\frac{\ell+5}4(\sum_p\bar\rho_p\log p+\log\pi)+$ [a nonnegative quantity
above $2,3,5,\ell$], where $\bar\rho_p=\sum_{v\mid p}w_v\rho_v$ is the weighted average of the outer radii of the
log-shells above $p$. So
Corollary~3.12 becomes $(k-1)a\le B$, the height inequality itself with the radii of the log-shells as margin (``in the
reading of IV, Corollary~3.12 carries the whole wall'' of \cite{RF}).

\tagH{} In terms of computer science, Corollary~3.12 is a relation (a relational property) between the outputs of two
programs (the $q$-side and the $\Theta$-side) \emph{over all inputs}. On the six completely factored data the comparison
holds because the wall is small, 5--13\% of $C$, and the margin is large (Table~2; D2357 falls short for its counted
places, and its values are conditional since the places above its unfactored cofactor are not counted). As a test, nothing has been tried. Turning this into a
proof is supposed to be the task of the multiradial algorithm of Theorem~3.11 (\S\S5, 6).

\paragraph{(vi) What the oracle \orc{step-xi} hides.} \tagP{} In the mockup, the only place where ``for every initial
$\Theta$-datum'' can enter is the oracle \orc{step-xi}. The implemented part (A)--(G) computes the same function of the
local invariants on actual data and on the data of $q\to q^N$. The margin $B$ of the hull does not grow with
$|\log\qq|$ (it moves slightly through $\gcd(15\ell,\mathrm{ord}\,q)$; for D211, $C$ takes two values,
\scanCvalues). Hence on the family $q\to q^N$ the oracle \orc{step-xi} answers ``no'' from $N=3$ (from $N=4$ for D2).

\tagH{} If a genuine Step~(xi) holds for all genuine data, it must be using something that the implemented part does not
have: a property that couples the size of the possible images with $|\log\qq|$, or facts that distinguish genuine data
from arbitrary local data. This is the same as lana-agents/iut's ``no other relation coupling $V$ with $a$'' and ``facts
about $D$ beyond them'' (\S6.3), seen in the program. When the oracles are replaced by implementations, this is what
remains last.

\paragraph{(vii) Cyclotomic rigidity fixes the scale.} \tagP{} Without step (C), the $j^2$-structure of the Kummer classes
($\kappa(\theta_j)=j^2\kappa(\theta_1)$) is still seen, but the size of $q$ is not: the values become multiples of the
unit chosen by the theater, and the output swings on the order of $10^{18}$. With step (C), the output does not depend on
the choice of coordinates.

\tagH{} Within the finite model, cyclotomic rigidity works as what fixes the scale when \'etale-like classes are turned
into numbers. That it is obtained from the structure of a group alone, the commutator in the theta group, agrees with
the remark of the paper on the \'etale theta function (quoted in \S9.2). This is, however, a finite model of the
algebraic core, not the cyclotomic rigidity of IUT itself (\S11).

\paragraph{(viii) Multiradiality as non-interference.} \tagH{} The description in the progress report (a common container
built from the inputs common to the two sides alone; \S6.1) is, in terms of computer science, a requirement of
non-interference. In the mockup this is enforced for the implemented part by the structure of the headers and checked
by changing the coordinates of the theaters. The oracles, on the other hand, read the rings of the theaters. Replacing
an oracle by an implementation means removing that reading of the ring while keeping the output; the separation by
types then becomes a tool for checking mechanically that the replacement is non-interfering.

\section{Limitations}
\begin{itemize}
\item \textbf{Not all readings are covered.} The 36 readings are the combinations of the axes found so far by the author
and by lana-agents/iut. We do not claim that they exhaust the readings of the sources. Moreover, the C program and the C++ mockup
implement only one reading, that of IV (the 36 readings are in the Python programs of \cite{RF}).
\item \textbf{A computation is not a proof.} That both sides can be computed, and that the inequality holds on the six
completely factored data, shows nothing about the correctness of IUT or about the Szpiro conjecture.
\item \textbf{The mockup is not the construction of IUT.} The theaters, the $\Theta$-link and the theta group are finite
models, not the objects of IUT (Frobenioids, tempered fundamental groups, the log-theta-lattice). The cyclotomic rigidity
is also only a finite model of the algebraic core, the commutator in a theta group. (Ind1), (Ind2), (Ind3) are the
formulas of the reading of IV, not the indeterminacies of Theorem~3.11 themselves.
\item \textbf{The values of the oracles are reverse-computed.} In particular, the line ``Corollary 3.12 (from the oracle
step-xi)'' of the output of the mockup puts the conclusion in and takes the conclusion out. The mockup only shows where
the algorithm has to enter and what is needed there.
\item \textbf{Assumptions of the finite level.} The Kummer classes are taken at level $N=2^{61}-1$, and the program checks
that $N$ is prime to $p(q_v-1)$. That the valuations of the theta values are integers (that $\qq$ lies in $K_v$) holds
because $\mathrm{ord}\,q$ is even at the bad places of these Legendre data; the mockup checks this too and stops if it
fails.
\item \textbf{The two programs and the mockup were all written on the author's side} (with generative AI). The
``classical computation'' was written from the statement of Theorem~1.10 without IUT, but it is not an implementation by
an independent third party. In the C programs and the mockup the extraction of the local data is also shared (in the
Python programs the two were extracted independently and agreed).
\item \textbf{The local data involve choices and assumptions.} $\Vbad$ is the subset of that of the catalogue without the
places above 3 and 5. The places above $2,3,5,\ell$ are set to 0 in the hull and bounded by Step~(ii) and (R4) in the
upper bound. The ramification index of $K$ at bad places is assumed by the formula of the catalogue except above $p$.
D2357 has an unfactored cofactor of 1564 digits whose places are not counted.
\item \textbf{There are seven data; the six completely factored ones have a large margin, and the values for D2357 are
conditional} (the places above its unfactored cofactor are not counted; that it is an initial $\Theta$-datum is proved
in \cite[Theorem~2.1]{Fourteen}). The family $q\to q^N$ is an adversarial input, but in
general it does not consist of initial $\Theta$-data of actual curves. None of the tests tries the claims of IUT.
\item \textbf{The scope of novelty is narrow.} The view of two computations is in Mochizuki and in the report of LANA;
the explicit computation of the IUT side is printed in IUT~IV and redone, with its constant, by Dupuy--Hilado and Zhou;
the objects and processes of IUT are C data types and a Lean port, with a numerical form of Step~(xi) and the same
constant, in motizukipdf; for one reading the equivalence with the height inequality is proved in Lean in
Takkun-kohinata/IUT\_LEAN; and the division into a finite layer and a structural layer is in Maeki (\S1.3). What is new
here is limited to placing the classical computation and the computation of IUT as programs on the same input, the
mockup with oracles separated by types, and the computer-science reading and position. The search for prior work was
made on 5 November 2026; version 1.0.0 did not cite motizukipdf or Takkun-kohinata/IUT\_LEAN, and the search may have
missed other things.
\item \textbf{``The essence'' and ``not understood essentially'' are the position of this note} (\tagH). Mochizuki also
speaks of the content of IUT in terms of multiradiality, anabelian geometry, Kummer theory and so on, and our reading of
his words is our own. There are opposite positions. Joshi holds that the left end also changes with the reading.
Scholze--Stix hold that the proof of Corollary~3.12 has a gap, that if the identifications of the copies of the real
numbers are taken consistently the inequality becomes empty, and that blurring by a factor of at least
$\mathrm{O}(\ell^2)$ makes it useless \cite[\S2.2]{SS}.
\item \textbf{The computer-science viewpoint remains a proposal.} The frameworks of equivalence, information flow and
relational verification are not formally applied in this note. The check of non-interference is also a test of one
hidden choice (the coordinates on $\Delta_\Theta$).
\item \textbf{No line-by-line correspondence can be asked for.} As Mochizuki's caveat (\S7) says, explicitly tracing the
correspondence between the parts of the two computations is something that cannot be done at all. Our comparison stays
at the shape of the whole computation.
\item \textbf{Sizes such as line counts mean nothing.} The length of a program depends on how it is written and does not
measure the difficulty of the computation.
\item \textbf{lana-agents/iut is moving.} We checked commit \texttt{553076a} (3 November 2026, the latest as of 5
November). The plan (\lean{Iut/\allowbreak{}Multiradial/\allowbreak{}ROADMAP.md}) is a draft of 14 October 2026.
\end{itemize}

\section{Future work}
\tagH{} The future task is to implement the oracles as actual algorithms, or as approximations to them. In the order of
Table~5:
\begin{enumerate}
\item \orc{theta-group} and \orc{kummer}: build the mono-theta environments at a finite level and the Kummer classes of
the theta values from finite quotients of tempered fundamental groups, without reverse computation from the ring of the
theater. lana-agents/iut has the Heisenberg group, the Kummer classes of the theta function, and parts of the discrete
rigidity.
\item \orc{gab}: read the type of the abstract group from a finite presentation of the group, not from the field. The
recovery of the local invariants is within known results.
\item \orc{ind3}: add the iteration of log-links and the log-Kummer correspondence, and make the containment of (Ind3) a
computation rather than an assumption.
\item \orc{global}: obtain the global data from the global anabelian reconstruction. By the plan of lana-agents/iut,
$p$-adic Hodge theory is needed here.
\item \orc{step-xi}: last, Step~(xi). After the replacements above, ask whether the output contains the $q$-pilot of
$\ddagger$, without reverse computation from the conclusion. As \S10 (vi) says, what is needed here is a property of ABC
strength.
\end{enumerate}

\tagH{} As ``approximation algorithms'' one may think of truncating the infinite constructions to finite ones: a finite
level $N$, finite quotients of profinite groups, finitely many columns of the log-theta-lattice. The questions to ask
are then whether the output converges as the truncation is refined, and whether the comparison corresponding to
Step~(xi) is stable under the truncation. In C or C++, replacing infinite profinite groups by finite objects is a
premise, and which finite models are faithful enough is itself a question.

\tagH{} On the side of tools one can use relational verification, which handles relations over all inputs; type systems
of information flow, which enforce multiradiality by types; and formalization in a proof assistant (the plan of
lana-agents/iut estimates the formalization of the whole of Theorem~3.11 in Lean at 0.7 to 1.5 million lines). Each of
them works as a check when the oracles are removed one by one.

\section{Reproducing}
The folder \texttt{two-computations/} of the bundle contains the C programs (\texttt{c/}), the C++ mockup
(\texttt{cpp/}), the seven local data (\texttt{data/}), the scripts (\texttt{scripts/}) and the outputs
(\texttt{outputs/}); its \texttt{README.md} lists the commands; \lean{scripts/run\_all.sh} rebuilds the programs, rewrites \texttt{outputs/}
and runs all the checks (about eight minutes), and \lean{paper/make\_generated.py} writes the tables and the code
excerpts of this note from the outputs and the sources. The C programs compile with
\texttt{gcc -O2 -Wall -Wextra -std=c11} without warnings; D2357 ($\ell=157$, long tuples counted by linearity) takes
about 10 seconds and the others less than a second. The mockup compiles with
\texttt{g++ -std=c++17 -O2 -Wall -Wextra -pedantic} without warnings (\texttt{make}), and \texttt{make check} checks the
information flow and the non-interference. \lean{check\_c.py} compares the C programs with the Python programs of
\cite{RF} (119 items; for D2357 the Python programs add a lower bound 261.75 of the container for the places above the
unfactored cofactor, which the C program does not count, so the comparison is made on the counted places), and
\lean{check\_cpp.py} compares the mockup with the C programs (\cppitems{} items) and runs the checks of
non-interference and the scan of $q\to q^N$. The quotations of Mochizuki's documents, of the report of LANA, of
Dupuy--Hilado, Zhou, Maeki and Joshi, and of lana-agents/iut were taken verbatim from the sources and checked by
machine, and the mathematical symbols in them on the page images. The quotations added in version 1.1.0 (of IUT~IV, of
Geiger and of Takkun-kohinata/IUT\_LEAN) were checked against the source texts, and those of IUT~IV and Geiger also on
the page images.

\section*{Use of generative AI}
The question studied in this note --- whether placing the computation of IUT and the classical computation side by side
as programs lets one approach the content of IUT from the viewpoint of computer science --- and the approach of regarding
the multiradial algorithm as the essence and placing it in a mockup as oracles arose in the author's discussions with
generative AI tools; the author judged them important and decided to make them the subject of this note. Its ultimate
aim, to understand exactly the algorithm by which IUT computes the height inequality, is the author's research
direction. Generative AI tools (large language models) were used extensively, under the author's direction, in the
computations, in writing and checking the programs, and in the writing. The author takes full responsibility for the
contents of this note.

\appendix
\section{Output of the mockup for D211}
\VerbatimInput{../outputs/mockup_D211.txt}
\VerbatimInput{../outputs/selftest_D211.txt}
\VerbatimInput{../outputs/scan_D211.txt}

\begin{thebibliography}{AbsTopIII}\raggedright
\bibitem[AbsTopIII]{AbsTopIII}
Shinichi Mochizuki, \emph{Topics in absolute anabelian geometry {III}: global reconstruction algorithms}, J. Math. Sci.
Univ. Tokyo \textbf{22} (2015), 939--1156.

\bibitem[Alien]{Alien}
Shinichi Mochizuki, \emph{The mathematics of mutually alien copies: from {G}aussian integrals to inter-universal
{T}eichm\"uller theory}, in: Inter-universal Teichm\"uller Theory Summit 2016, RIMS K\^oky\^uroku Bessatsu \textbf{B84}
(2021), 23--192.

\bibitem[Cat]{Cat}
Toshihisa Urahigashi, \emph{Explicit initial $\Theta$-data and the two readings of Step (v) of IUT IV, Theorem 1.10:
a catalogue}, version 1.0.0, Zenodo, 2026, \href{https://doi.org/10.5281/zenodo.23870448}{doi:10.5281/zenodo.23870448}.

\bibitem[DH]{DH}
Taylor Dupuy and Anton Hilado, \emph{Probabilistic {S}zpiro, baby {S}zpiro, and explicit {S}zpiro from {M}ochizuki's
{C}orollary 3.12}, arXiv:2004.13108v2 [math.NT], 29 April 2020.

\bibitem[EssLgc]{EssLgc}
Shinichi Mochizuki, \emph{On the essential logical structure of inter-universal {T}eichm\"uller theory in terms of
logical {AND} ``$\wedge$''/logical {OR} ``$\vee$'' relations: report on the occasion of the publication of the four main
papers on inter-universal {T}eichm\"uller theory}, RIMS Preprint 1968, November 2022; version of March 2024.

\bibitem[EtTh]{EtTh}
Shinichi Mochizuki, \emph{The \'etale theta function and its {F}robenioid-theoretic manifestations}, Publ. Res. Inst.
Math. Sci. \textbf{45} (2009), no.~1, 227--349.

\bibitem[Fifth]{Fifth}
Toshihisa Urahigashi, \emph{Two constants in the numerical estimates of inter-universal {T}eichm\"uller theory {IV}:
$12/(\ell(\ell+1))$ and $12/\ell^2$}, fifth report, version 1.0.0, Zenodo, 2026,
\href{https://doi.org/10.5281/zenodo.22659796}{doi:10.5281/zenodo.22659796} (all versions).

\bibitem[First]{First}
Toshihisa Urahigashi, \emph{An explicit initial $\Theta$-datum with mixed reduction at a split prime}, first report,
version 0.1.7, Zenodo, 2026, \href{https://doi.org/10.5281/zenodo.22537342}{doi:10.5281/zenodo.22537342} (all
versions).

\bibitem[Form]{Form}
Shinichi Mochizuki, \emph{On the formalization of IUT: a preliminary progress report} (joint work in progress with
Y.~Hoshi, G.~Yamashita, Y.~Yang), slides, April 2026, 14 pp.,
\url{https://www.kurims.kyoto-u.ac.jp/~motizuki/Formalization\%20of\%20IUT\%20(2026-04\%20version).pdf}.

\bibitem[Fourteen]{Fourteen}
Toshihisa Urahigashi, \emph{Fourteen initial $\Theta$-data on one curve over $\Q(i)$}, Zenodo, 2026,
\href{https://doi.org/10.5281/zenodo.23990716}{doi:10.5281/zenodo.23990716} (all versions).

\bibitem[Fourth]{Fourth}
Toshihisa Urahigashi, \emph{Split bad primes and the tuple structure of inter-universal {T}eichm\"uller theory}, fourth
report, version 1.0.4, Zenodo, 2026, \href{https://doi.org/10.5281/zenodo.22648944}{doi:10.5281/zenodo.22648944} (all
versions).

\bibitem[Gap]{Gap}
Toshihisa Urahigashi, \emph{The gap between the two walls is positive: a closed form for the (Ind2)-terms of
lana-agents/iut}, version 1.0.0, Zenodo, 2026, \href{https://doi.org/10.5281/zenodo.23954194}{doi:10.5281/zenodo.23954194}.

\bibitem[Ge2]{Ge2}
Lukas Geiger, \emph{The {B}ridge-{C}ontract {C}riterion: a form audit for transport arguments, with an instantiation in
the abc literature} (The abc Gap Series, Part~2), version 0.2, Zenodo, 16 August 2026,
\href{https://doi.org/10.5281/zenodo.21960477}{doi:10.5281/zenodo.21960477}; all versions
\href{https://doi.org/10.5281/zenodo.21960476}{doi:10.5281/zenodo.21960476}.

\bibitem[Hi]{Hi}
Joaquim Reizi Higuchi, \emph{A formal negative result for the passage from {S}tep (xi-e) to {S}tep (xi-f) in {IUT}
{III}}, preprint, 31 March 2026, \href{https://doi.org/10.13140/RG.2.2.10989.76006}{doi:10.13140/RG.2.2.10989.76006}.

\bibitem[II]{II}
Toshihisa Urahigashi, \emph{Reading the IUT dispute through indexed identifications: an educational model in
undergraduate mathematics}, version 2.0.0, Zenodo, 2026,
\href{https://doi.org/10.5281/zenodo.23933108}{doi:10.5281/zenodo.23933108}.

\bibitem[IUT1]{IUT1}
Shinichi Mochizuki, \emph{Inter-universal {T}eichm\"uller theory {I}: construction of {H}odge theaters}, Publ. Res.
Inst. Math. Sci. \textbf{57} (2021), 3--207.

\bibitem[IUT3]{IUT3}
Shinichi Mochizuki, \emph{Inter-universal {T}eichm\"uller theory {III}: canonical splittings of the log-theta-lattice},
Publ. Res. Inst. Math. Sci. \textbf{57} (2021), 403--626.

\bibitem[IUT4]{IUT4}
Shinichi Mochizuki, \emph{Inter-universal {T}eichm\"uller theory {IV}: log-volume computations and set-theoretic
foundations}, Publ. Res. Inst. Math. Sci. \textbf{57} (2021), 627--723.

\bibitem[Joshi]{Joshi}
Kirti Joshi, \emph{Comments on the LANA Project Report of Kato et al.}, 30 July 2026,
\url{https://sites.arizona.edu/kirti-joshi/files/2026/07/Comments-on-the-LANA-Project-Report-of-Kato-et-al.pdf}.

\bibitem[lana]{lana}
lana-agents, \emph{iut}: a Lean~4 formalization of inter-universal {T}eichm\"uller theory, GitHub repository,
\url{https://github.com/lana-agents/iut}, commit \texttt{553076a} (3 November 2026). Files cited:
\lean{Iut/\allowbreak{}WallEquiv/\allowbreak{}NonDerivable.lean}, \lean{Iut/\allowbreak{}WallEquiv/\allowbreak{}README.md}, \lean{Iut/\allowbreak{}MainTheorem.lean},
\lean{Iut/\allowbreak{}EtaleTheta/\allowbreak{}Heisenberg.lean}, \lean{Iut/\allowbreak{}Multiradial/\allowbreak{}ROADMAP.md} (a draft of 14 October 2026).

\bibitem[LANA]{LANA}
Project LANA, \emph{Project LANA interim report on IUT theory}, 20 July 2026,
\url{https://ncatlab.org/nlab/files/LANAProject-Report-July2026.pdf}.

\bibitem[Maeki]{Maeki}
Hidemitsu Maeki, \emph{IUT-Explicit-ADIC}, Zenodo, 22 November 2025,
\href{https://doi.org/10.5281/zenodo.17677052}{doi:10.5281/zenodo.17677052}.

\bibitem[Pj]{Pj}
Paul Pajo, \emph{Scale-free common-container comparisons do not entail numerical inequalities: a minimum-viable no-go
theorem motivated by the {IUT} {III} {T}heorem 3.11-to-{C}orollary 3.12 bridge}, preliminary research note,
26 April 2026, \href{https://doi.org/10.13140/RG.2.2.33620.16005}{doi:10.13140/RG.2.2.33620.16005}.

\bibitem[Richards]{Richards}
David Richards, \emph{motizukipdf---{M}ochizuki's papers, plus a toolkit for the data types of {IUT}}, GitHub
repository \url{https://github.com/oneman/motizukipdf}, created on 27 October 2026, consulted at commit \texttt{2665db1}
(30 October 2026); Lean port \lean{lean/IUT/LeanForm.lean}.

\bibitem[RF]{RF}
Toshihisa Urahigashi, \emph{No reading avoids the height inequality: 36 readings of Corollary 3.12 and the conservation
of the height load}, version 1.0.0, Zenodo, 2026,
\href{https://doi.org/10.5281/zenodo.23965952}{doi:10.5281/zenodo.23965952}.

\bibitem[Second]{Second}
Toshihisa Urahigashi, \emph{An initial $\Theta$-datum over an imaginary quadratic field, and numerical values of the
local quantities of inter-universal {T}eichm\"uller theory {IV}}, second report, version 1.0.5, Zenodo, 2026,
\href{https://doi.org/10.5281/zenodo.22601333}{doi:10.5281/zenodo.22601333} (all versions).

\bibitem[SL]{SL}
Toshihisa Urahigashi, \emph{Segment or loop? A homotopy-type-theoretic overview of the dispute over Corollary 3.12 of
inter-universal {T}eichm\"uller theory}, version 1.1.0, Zenodo, 2026,
\href{https://doi.org/10.5281/zenodo.23933109}{doi:10.5281/zenodo.23933109}.

\bibitem[SS]{SS}
Peter Scholze and Jakob Stix, \emph{Why abc is still a conjecture}, 2018.

\bibitem[Stepv]{Stepv}
Toshihisa Urahigashi, \emph{A local non-containment in Step (v) of IUT IV, Theorem 1.10}, version 3.0.0, Zenodo, 2026,
\href{https://doi.org/10.5281/zenodo.23852752}{doi:10.5281/zenodo.23852752}.

\bibitem[Third]{Third}
Toshihisa Urahigashi, \emph{The cost of each distinction in the local computations of inter-universal {T}eichm\"uller
theory {IV}}, third report, version 1.0.2, Zenodo, 2026,
\href{https://doi.org/10.5281/zenodo.22646136}{doi:10.5281/zenodo.22646136} (all versions).

\bibitem[TK]{TK}
Takkun-kohinata, \emph{IUT\_LEAN}, Lean repository, \url{https://github.com/Takkun-kohinata/IUT_LEAN}, commit
\texttt{9d56c46} (15 July 2026); files \lean{IutLean/Multiradial/ThetaHullDichotomy.lean} and
\lean{REQUEST\_theta\_hull\_derivation.md}.

\bibitem[W311]{W311}
Toshihisa Urahigashi, \emph{Where Theorem 3.11 sits: the log-volume content of the results cited in the proof of IUT
III, Corollary 3.12, does not imply its conclusion}, Zenodo, 2026,
\href{https://doi.org/10.5281/zenodo.24028832}{doi:10.5281/zenodo.24028832} (all versions).

\bibitem[WT]{WT}
Toshihisa Urahigashi, \emph{A programmer's walkthrough of inter-universal {T}eichm\"uller theory: running the
multiradial algorithm on explicit initial $\Theta$-data}, version 1.1.0, Zenodo, 2026,
\href{https://doi.org/10.5281/zenodo.23955994}{doi:10.5281/zenodo.23955994}.

\bibitem[Zhou]{Zhou}
Zhong-Peng Zhou, \emph{The inter-universal {T}eichm\"uller theory and new {D}iophantine results over the rational
numbers.~I}, arXiv:2503.14510v1 [math.NT], 8 March 2025; code at \url{https://github.com/zhongpengzhou/IUT-Q-I-code}.

\bibitem[ZhouII]{ZhouII}
Zhong-Peng Zhou, \emph{The inter-universal {T}eichm\"uller theory and new {D}iophantine results over rational
numbers.~{II}}, arXiv:2510.05448v1 [math.NT], 6 October 2025; code at
\url{https://github.com/zhongpengzhou/IUT-Q-II-code}.
\end{thebibliography}

\end{document}
